\section{Related Work}
\label{sec:related}

\paragraph{Aerial person detection for search and rescue.}
HERIDAL~\cite{bozic2019heridal} and SARD~\cite{sambolek2021sard} are the
two datasets this task is normally validated on. They differ in capture
regime, HERIDAL at high altitude over wilderness terrain with very small
subjects and SARD at lower altitude with larger relative object size,
and together they cover most of what the literature reports. Both were
collected in clear daytime conditions, and both are evaluated that way.
VisDrone~\cite{zhu2022visdrone} is broader in scene content and includes
genuine night frames, but it is urban and artificially lit, and person
annotations there describe a different operating context than wilderness
search. We use it only as an independent check, never for training,
calibration or threshold selection.

\paragraph{Corruption robustness benchmarks.}
ImageNet-C~\cite{hendrycks2019imagenetc} established the pattern this
work follows: hold the dataset fixed, apply a fixed corruption ladder,
and report accuracy per condition rather than as a single number.
Michaelis \etal~\cite{michaelis2019benchmarking} carried the protocol to
detection and showed stylisation-based augmentation recovers a
substantial part of the loss, which is the independent precedent our
mitigation arm draws on. These benchmarks are built from generic image
operations chosen for breadth. That is the right choice for a general
recognition benchmark and the wrong one here, because a responder
selecting a system for a night flood mission needs a number attached to
a named condition, not to a filter. Foggy Cityscapes~\cite{sakaridis2018foggy}
is the closest precedent for the alternative: it synthesises fog from
the Koschmieder model on real imagery with depth, and demonstrates that
a physically derived corruption applied to real scenes supports
controlled attribution. AegisBench extends that stance from one
condition in one domain to nine conditions across four disaster families
in aerial Earth observation, and adds a declared per-corruption
calibration statistic so that severity claims are falsifiable rather
than stipulated.

\paragraph{Optical models of degraded imaging.}
The corruption implementations are not new physics. Atmospheric
scattering follows Koschmieder~\cite{koschmieder1924theorie} in the form
standardised for vision by Narasimhan and
Nayar~\cite{narasimhan2002vision}; rain streaks follow the additive
directional model of Garg and Nayar~\cite{garg2007rain}; low light
follows the linear-domain photometric account of photon-limited capture
used by Chen \etal~\cite{chen2018dark}, including the signal-dependent
noise that dominates as light falls. The contribution is not the models
but their use as a calibrated, reproducible evaluation axis for a
safety-of-life aerial task, together with the parameter ranges,
calibration statistics and seeding discipline that make a severity mean
the same thing across capture regimes.

\paragraph{What is missing.}
Across these lines of work, robustness for aerial search and rescue is
reported either under clear conditions or under generic distortions, and
almost always at a tuned operating point. None of the three gaps is
exotic: what is absent is a benchmark that names the disaster condition,
calibrates its severity, and scores every architecture through one
evaluator at a threshold a fielded system could actually have chosen in
advance. That combination is what AegisBench supplies, and it is what
makes the low-light result in Section~\ref{sec:results} interpretable as
an operational finding rather than a tuning artefact.
